%% Figure: the pumping lemma (Lemma 3.3) and the proof of the speed-1/3 barrier (Theorem 3.4).
%% Schematic. In (a) the relative cell (tau,kappa) of R_m is [m-kappa-1,m-kappa] x [-tau-3,-tau-2].
\begin{tikzpicture}
  % (a) R_n and R_n' with the same values on the cone of the right end of R_n (n = 8, n' = 12)
  \begin{scope}[x=2.7mm,y=2.7mm]
    \foreach \m/\dx in {8/0,12/12} {
      \begin{scope}[shift={(\dx,0)}]
        \foreach \tau in {-2,...,11} {\foreach \kap in {0,...,11} {%
          \pgfmathtruncatemacro{\cls}{((\kap<=\tau)&&(\tau<=\m-1)) ? ((\tau==\m-1) ? 1 : ((\tau+\kap<=7) ? 2 : 3))
            : ((((\tau==\kap-1)||((\tau==\kap-2)&&(\kap>=1)))&&(\kap<=\m-1)) ? ((\kap<=4) ? 4 : 5) : 0)}%
          \ifnum\cls>0
            \def\cc{\ifcase\cls\or stF\or rgRefl2\or rgRefl\or rgCone\or rgInput\fi}%
            \filldraw[fill=\cc,draw=grid!45,line width=0.15pt] (\m-\kap-1,-\tau-2) rectangle (\m-\kap,-\tau-3);
          \fi}}
        \draw[ink,line width=1.2pt] (\m,0.1) -- (\m,-\m-2.1);
        \draw[ink,line width=1.2pt] (0,-\m+0.9) -- (0,-\m-2.1);
        \node[star,star points=5,fill=white,draw=ink,line width=0.4pt,inner sep=0pt,minimum size=0.24cm]
          at (\m-0.5,-9.5) {};
      \end{scope}}
    \node[font=\small,anchor=south] at (4,0.4) {$R_n$, $n=8$};
    \node[font=\small,anchor=south] at (20.5,0.4) {$R_{n'}$, $n'=12$};
    \draw[-{Stealth[length=1.8mm]},ink,line width=0.6pt] (7.6,-0.2) .. controls (9.6,2.7) and (14.4,2.7) .. (16.4,-0.2);
    \node[font=\scriptsize,anchor=south] at (12.0,2.25) {equal on the cone};
    \node[font=\scriptsize,anchor=north,align=center] at (4,-10.5) {$R_n(n-1,0)=\Fq$};
    \node[font=\scriptsize,anchor=north,align=center] at (18,-14.5) {then also $R_{n'}(n-1,0)=\Fq$,\\
      although $n-1<n'-1$};
  \end{scope}
  % (b) the half-line: the cones of the lengths n and n+P on the depth rows
  \begin{scope}[shift={(9.2cm,0cm)},x=1.6mm,y=1.6mm]
    \fill[gNone] (0,0) -- (36,-36) -- (0,-36) -- cycle;
    \foreach \j in {4,8,12} \draw[grid,line width=0.45pt] (0,-\j) -- (36-\j,-36);
    \draw[ink,densely dashed,line width=0.7pt] (0,0) -- (12,-36);
    \draw[rgCone,line width=2.4pt] (14,-14) -- (7,-21);
    \draw[rgCone,line width=2.4pt] (20,-20) -- (10,-30);
    \foreach \j in {4,8,12} {
      \fill[ink] (14-\j/2,-14-\j/2) circle (0.45);
      \fill[ink] (20-\j/2,-20-\j/2) circle (0.45);}
    \draw[-{Stealth[length=1.6mm]},stF,line width=0.8pt] (10.8,-18.8) -- (15.3,-23.3);
    \node[font=\scriptsize,anchor=south west,inner sep=1pt] at (13.0,-20.6) {$+P$};
    \draw[ink,line width=1.2pt] (0,0) -- (0,-36);
    \draw[ink,line width=0.7pt] (0,0) -- (36,-36);
    \node[font=\scriptsize,rotate=-45,anchor=south] at (6.5,-6.5) {front};
    \node[font=\scriptsize,anchor=north] at (12,-36.2) {$i=t/3$};
    \node[font=\scriptsize,anchor=north] at (29,-36.2) {depth rows};
    \node[font=\scriptsize,anchor=east,align=right] at (-0.8,-15.0) {cone of\\ the length $n$};
    \draw[-{Stealth[length=1.4mm]},ink,line width=0.45pt] (-0.7,-15.6) -- (8.6,-19.2);
    \node[font=\scriptsize,anchor=west,align=left] at (26.0,-19.0) {cone of the\\ length $n+P$};
    \draw[-{Stealth[length=1.4mm]},ink,line width=0.45pt] (25.9,-19.6) -- (17.2,-22.6);
  \end{scope}
  \node[font=\small,anchor=south west] at (-0.3,1.25) {(a) the pumping lemma};
  \node[font=\small,anchor=south west] at (8.2,1.25) {(b) the proof of \cref{thm:third}};
\end{tikzpicture}
